\section{iPRA, finite codes, and the meaning of identity}
\label{sec:arithmetic}

\subsection{The object theory}

We use first-order intuitionistic logic with equality over the natural
numbers. The language has $0$, successor, and a symbol for every primitive
recursive definition. Each symbol has its defining equations. The
nonlogical axioms also include successor injectivity, $Sx\ne0$, and
induction for quantifier-free formulas, with parameters:
\[
 A(0,\boldsymbol y)\land
 \forall x\bigl(A(x,\boldsymbol y)\to A(Sx,\boldsymbol y)\bigr)
 \ \longrightarrow\ \forall x\,A(x,\boldsymbol y).
\]
The formula $A$ must be syntactically quantifier-free. Classical excluded
middle for arbitrary formulas is not an axiom. Equality and order on
natural numbers, integers encoded by natural numbers, and rational numbers
are decidable by primitive recursive comparisons.

A useful derived rule is induction for each fixed bounded formula whose
atomic relations are primitive recursive. Its characteristic function is
constructed by finite Boolean folds. A bounded existential is a finite
disjunction and a bounded universal is a finite conjunction. The resulting
formula is internally equivalent to a quantifier-free equation
$\chi(\boldsymbol x)=1$, to which the original induction scheme applies.
This observation does not grant induction for an unbounded existential
formula.

\subsection{Coding conventions}

Any explicit elementary coding with primitive recursive constructors and
projections works. For definiteness, use Cantor pairing
\[
 \langle a,b\rangle=\frac{(a+b)(a+b+1)}2+b.
\]
The division is exact; the inverse is obtained by bounded search below the
code. Encode an integer $z$ by $2z$ when $z\ge0$ and by $-2z-1$ when
$z<0$. Encode a rational as a pair consisting of an integer code and a
natural number $d$, interpreted as denominator $d+1$. Reduction to lowest
terms is optional. Exact rational comparisons cross-multiply positive
denominators.

Finite lists can be encoded by their length and a nested sequence of pairs.
Their projections, bounded sums, products, maxima, and bounded searches are
primitive recursive. An expression is a finite list of tagged nodes in
which every child index is smaller than its parent index. The final node
is the root. Tags distinguish rational constants, the variable, $\pi$,
$\log2$, addition, multiplication, sine, exponential, and absolute value.
The predicate $\WF(e)$ checks tags, arities, denominators, and child-index
bounds. A tree can be expanded from a directed acyclic graph if literal
expression syntax is desired. All conversions are primitive recursive.

A polynomial is a finite list of pairs $(c,\boldsymbol\alpha)$ with
$c\in\Z$ and $\boldsymbol\alpha\in\N^n$. Its value is
$\sum c\prod_i x_i^{\alpha_i}$. Zero coefficients and repeated monomials
are permitted; collecting like terms is not required for any proof.
The degree and coefficient-norm bounds used below are computed from this
finite list. A dummy variable handles a polynomial with no variables.

These conventions specify ordinary data, rather than adding new sorts to
iPRA. Symbols such as $q\in\Q$ or $\boldsymbol m\in\N^n$ abbreviate
natural-number codes and their primitive recursive decoding operations.

\subsection{Rational approximation and finite evidence}

\begin{lemma}[Uniform approximation]\label{lem:approx}
There is a primitive recursive rational-valued function $\App(e,q,k)$
such that, for every well-formed expression $e$ and rational $q$,
\begin{equation}\label{eq:approx-error}
 \abs{\App(e,q,k)-\llbracket e\rrbracket(q)}\le2^{-k}.
\end{equation}
The elementary approximation laws needed in the proof have iPRA
derivations after replacing real-valued notation by their finite rational
error statements.
\end{lemma}

The brackets denote the ordinary real interpretation for readability.
Appendix~\ref{app:primitives} gives concrete Taylor cutoffs, approximations
to the constants, and a global primitive recursive precision bound.
In particular, this lemma is not justified by an alleged primitive
recursive interpreter for arbitrary primitive recursive programs. Its
syntax has a fixed, small set of elementary operations.

\begin{definition}[Strict certificate and identity]
For a valid expression code $e$, put
\begin{align}
 \NZ(e,q,k)&:\Longleftrightarrow
       \abs{\App(e,q,k)}>2^{-k},\label{eq:nz}\\
 W(e)&:\Longleftrightarrow\exists q\,\exists k\;\NZ(e,q,k),\label{eq:w}\\
 \Id(e)&:\Longleftrightarrow
       \forall q\,\forall k\;\abs{\App(e,q,k)}\le2^{-k}.
       \label{eq:id}
\end{align}
Rational input codes range over a surjective enumeration of $\Q$.
\end{definition}

The relation $\NZ$ is a decidable primitive recursive relation. Thus $W$
is an existential finite-certificate statement and $\Id$ is a universal
arithmetical statement. Intuitionistic logic proves
\begin{equation}\label{eq:id-neg-w}
 \Id(e)\longleftrightarrow\neg W(e).
\end{equation}
For the right-to-left direction, fix $q,k$. The rational comparison is
decidable, so exclusion of its strict alternative proves the displayed
nonstrict inequality. No unbounded logical principle is involved.

\begin{proposition}[Agreement with real identity]\label{prop:semantics}
In the standard interpretation, $\Id(e)$ holds exactly when
$\llbracket e\rrbracket(x)=0$ for every real $x$.
Moreover, a specified nonzero value with a positive separation from zero
can be converted into a strict rational certificate using a local
continuity modulus.
\end{proposition}
\begin{proof}
If the function vanishes everywhere, \eqref{eq:approx-error} immediately
gives \eqref{eq:id}. Conversely, fix a rational $q$. If \eqref{eq:id}
holds then
\[
 \abs{\llbracket e\rrbracket(q)}
 \le \abs{\App(e,q,k)}+2^{-k}\le2^{1-k}
\]
for every $k$. This is precisely equality to zero in rational-approximation
semantics.

On every bounded interval, an expression has a computable rational
magnitude bound and a computable Lipschitz bound. Addition, multiplication,
sine, exponential on a bounded interval, and absolute value preserve such
bounds. Induction on the finite syntax therefore gives an explicit local
modulus. Approximating any real $x$ by rationals within that modulus shows
that its value is zero too. This does not require a choice function for all
real numbers: a supplied Cauchy name is evaluated at a specified precision.

For the positive assertion, suppose a supplied point has value of absolute
size at least $\delta>0$, where $\delta$ is rational. Approximate the point
closely enough that the function value changes by less than $\delta/4$.
Then approximate the value at this rational point with error
$2^{-k}<\delta/4$. Its rational approximation has absolute value greater
than $\delta/2>2^{-k}$, furnishing \eqref{eq:nz}.
\end{proof}

\begin{remark}[A logical distinction that must be preserved]
From \eqref{eq:id-neg-w}, intuitionistically $\neg\Id(e)$ implies
$\neg\neg W(e)$. It does \emph{not} in general imply $W(e)$.
Nowhere below do we replace a negated universal statement by a positive
counterexample. Our positive polynomial-to-expression equivalence is
proved directly with witness transformations, and only then negated.
\end{remark}

\subsection{Finite computations}

Fix a standard effective enumeration of deterministic machines whose
terminating instruction emits a Boolean output. This loses no halting
generality: simulate any ordinary program and emit $0$ when it halts.
Intermediate computations may manipulate arbitrary finite data. A
configuration records the state, a finite tape segment, and the head
position, or equivalently the data of a register-machine configuration.
A certificate records a finite sequence of configurations, its initial
input, each legal transition, and its final output. Write
$\Run(d,x,s,b)$ when $s$ is such a valid terminating computation of program
$d$ on input $x$ with output bit $b$, encoded by $0$ or $1$.
Verification is bounded by the length
and entries of the supplied certificate and is primitive recursive.

The fixed halting formula is
\begin{equation}\label{eq:halt-def}
 \Halt(e):\Longleftrightarrow
 \exists s\,\exists b\;\Run(e,e,s,b).
\end{equation}
Thus halting is on the machine's own input code, not on a blank tape.
One can encode configurations by a standard bounded arithmetic sequence
coding (for example the remainder coding used for computation tableaux).
Initiality, local transitions, and finality then have bounded arithmetic
definitions, and \eqref{eq:halt-def} is a specific $\Sigma_1$ formula.
If an exponential coding is used, its finitely many auxiliary computations
are included in the existential certificate. This is the fixed formula to
which formalized MRDP will be applied.

Two elementary internal facts are used later. Valid terminating
computations of a deterministic program on the same input have equal
outputs. Also, a finite program that calls another program and then
branches has primitive recursive transformations between the corresponding
finite run certificates. Both facts are proved by bounded induction on
the supplied computation sequences. There is no totality assumption for
the universal machine in these facts.
